\section{Motivation}

\subsection{Type Errors in Hindley--Milner Systems}

Hindley--Milner type inference operates by generating a set of type-equality
constraints from a program's syntax and solving them via Robinson unification~\cite{damas1982}.
When the program is ill-typed, this constraint set is unsatisfiable, and the compiler
must report an error location.
Standard implementations report the first unification failure encountered during
constraint solving, a location that is structurally determined by traversal order
rather than by any minimality criterion.
As a consequence, the reported location may not correspond to the actual source of the
type error, making error messages brittle and difficult to act on~\cite{pavlinovic2014}.

\subsection{The Algorithmic Significance}

A principled alternative is to identify a minimum-weight subset of constraints whose
removal restores satisfiability.
This is a combinatorial optimization problem.
Via the MCS--MUS hitting-set duality established by Marques-Silva et al.~\cite{marques2013},
this problem is directly related to finding a minimum-weight hitting set over the family
of Minimal Unsatisfiable Subsets of the constraint system.
Existing work addresses related localization problems using exact MaxSMT solvers~\cite{pavlinovic2014,kopinsky2024}
or MCS enumeration~\cite{fu2024}, both of which provide principled answers but can become
expensive as instance size grows.
The literature reviewed here does not provide polynomial-time approximation guarantees
or approximation-hardness results for this problem in the HM setting.

\subsection{Why Approximation Theory Matters Here}

Exact methods can be impractical for large programs, while heuristics without quality
guarantees offer no formal assurance on solution quality.
A polynomial-time approximation algorithm with a provable ratio would fill this gap by
trading optimality for scalability in a controlled, quantifiable way.
This paper takes a first empirical step: we implement a greedy approximation algorithm,
evaluate its approximation ratio against an exact baseline on synthetic HM-style
instances, and diagnose which structural features of those instances drive its behaviour.
